\nonstopmode{}
\documentclass[a4paper]{book}
\usepackage[times,inconsolata,hyper]{Rd}
\usepackage{makeidx}
\makeatletter\@ifl@t@r\fmtversion{2018/04/01}{}{\usepackage[utf8]{inputenc}}\makeatother
% \usepackage{graphicx} % @USE GRAPHICX@
\makeindex{}
\begin{document}
\chapter*{}
\begin{center}
{\textbf{\huge Package `statuser'}}
\par\bigskip{\large \today}
\end{center}
\ifthenelse{\boolean{Rd@use@hyper}}{\hypersetup{pdftitle = {statuser: Statistical Tools Designed for End Users}}}{}
\ifthenelse{\boolean{Rd@use@hyper}}{\hypersetup{pdfauthor = {Uri Simonsohn}}}{}
\begin{description}
\raggedright{}
\item[Type]\AsIs{Package}
\item[Title]\AsIs{Statistical Tools Designed for End Users}
\item[Version]\AsIs{0.2.1}
\item[Author]\AsIs{Uri Simonsohn [aut, cre]}
\item[Maintainer]\AsIs{Uri Simonsohn }\email{urisohn@gmail.com}\AsIs{}
\item[Description]\AsIs{The statistical tools in this package do one of four things:
1) Enhance basic statistical functions with more 
flexible inputs, smarter defaults, and richer, clearer, and ready-to-use 
output (e.g., t.test2())
2) Produce publication-ready commonly needed figures with one line of code (e.g., plot_cdf())
3) Implement novel analytical tools developed by the authors (e.g., twolines())
4) Deliver niche functions of high value to the authors that are not easily
available elsewhere (e.g., clear(), convert_to_sql(), resize_images()).}
\item[License]\AsIs{GPL-3}
\item[Encoding]\AsIs{UTF-8}
\item[RoxygenNote]\AsIs{7.3.3}
\item[Imports]\AsIs{mgcv, marginaleffects, rsvg, magick, sandwich, lmtest, utils}
\item[Suggests]\AsIs{testthat (>= 3.0.0), crayon, quantreg, estimatr, broom,
modelsummary}
\item[Config/testthat/edition]\AsIs{3}
\end{description}
\Rdcontents{Contents}
\HeaderA{statuser-package}{Stat Tools for End Users}{statuser.Rdash.package}
\aliasA{statuser}{statuser-package}{statuser}
\keyword{internal}{statuser-package}
%
\begin{Description}
Basic and custom statistical tools designed with end users in mind. 
Functions have optimized defaults, produce decluttered and informative output 
that is self-explanatory, and generate publication-ready results in 1 line of code.
\end{Description}
%
\begin{Section}{Basic Stats (improved)}

\begin{itemize}

\item{} \code{\LinkA{lm2}{lm2}}: like lm(), with robust SE and much more informative output
\item{} \code{\LinkA{t.test2}{t.test2}}: like t.test(), decluttered, and more informative output 
\item{} \code{\LinkA{table2}{table2}}: like table(), showing variable names, and with proportions \& chi2 built in
\item{} \code{\LinkA{desc\_var}{desc.Rul.var}}: Descriptive statistics for variables (optional, by group(s))

\end{itemize}

\end{Section}
%
\begin{Section}{Custom tools from papers by Uri Simonsohn}

\begin{itemize}

\item{} \code{\LinkA{twolines}{twolines}}: Two-lines test for U-shapes (Simonsohn 2018)
\item{} \code{\LinkA{interprobe}{interprobe}}: Probe and visualize nonlinear interactions (Simonsohn 2024; Montealegre \& Simonsohn 2026)

\end{itemize}

\end{Section}
%
\begin{Section}{Graphing}

\begin{itemize}

\item{} \code{\LinkA{scatter.gam}{scatter.gam}}: Makes scatter plot for x \& y, with fitted GAM line y=f(x)
\item{} \code{\LinkA{plot\_cdf}{plot.Rul.cdf}}: Plot empirical cumulative distribution functions (optional, by group)
\item{} \code{\LinkA{plot\_density}{plot.Rul.density}}: Plot density functions (optional, by group)
\item{} \code{\LinkA{plot\_freq}{plot.Rul.freq}}: Plot frequency of observed values (optional, by group)
\item{} \code{\LinkA{plot\_means}{plot.Rul.means}}: Barplot of means with confidence intervals and (optionally) tests
\item{} \code{\LinkA{plot\_gam}{plot.Rul.gam}}: Plot fitted GAM values for a focal predictor

\end{itemize}

\end{Section}
%
\begin{Section}{Formatting}

\begin{itemize}

\item{} \code{\LinkA{format\_pvalue}{format.Rul.pvalue}}: Format p-values for display
\item{} \code{\LinkA{message2}{message2}}: Print colored messages to console
\item{} \code{\LinkA{resize\_images}{resize.Rul.images}}: Resize images (SVG, PDF, EPS, JPG, PNG, etc.) to PNG with specified width

\end{itemize}

\end{Section}
%
\begin{Section}{Miscellaneous}

\begin{itemize}

\item{} \code{\LinkA{clear}{clear}}: Clear environment, console, and all graphics devices
\item{} \code{\LinkA{list2}{list2}}: Like list(), but unnamed objects are automatically named
\item{} \code{\LinkA{convert\_to\_sql}{convert.Rul.to.Rul.sql}}: Convert CSV files to SQL INSERT statements
\item{} \code{\LinkA{text2}{text2}}: like text() adding text-alignment and background color

\end{itemize}

\end{Section}
%
\begin{Author}
Uri Simonsohn \email{urisohn@gmail.com}
\end{Author}
%
\begin{References}
Simonsohn, U. (2018). Two lines: A valid alternative to the invalid testing of 
U-shaped relationships with quadratic regressions. \emph{Advances in Methods and 
Practices in Psychological Science}, 1(4), 538-555. \Rhref{https://doi.org/10.1177/2515245918805755}{doi:10.1177\slash{}2515245918805755}

Simonsohn, U. (2024). Interacting with curves: How to validly test and probe interactions in the real (nonlinear) world. \emph{Advances in Methods and Practices in Psychological Science}, 7(1), Article 25152459231207787. \Rhref{https://doi.org/10.1177/25152459231207787}{doi:10.1177\slash{}25152459231207787}

Montealegre, D., \& Simonsohn, U. (2026). \emph{Johnson Neyman 2.0} (working paper).
\end{References}
%
\begin{SeeAlso}
Useful links:
\begin{itemize}

\item{} \url{https://github.com/urisohn/statuser}

\end{itemize}

\end{SeeAlso}
\HeaderA{clear}{Clear Plot, Global Environment, and Console}{clear}
%
\begin{Description}
Clears plot, global environment, and console. On first use the user is
prompted to authorize clearing the environment to comply with CRAN rules.
\end{Description}
%
\begin{Usage}
\begin{verbatim}
clear()
\end{verbatim}
\end{Usage}
%
\begin{Details}
This function performs three cleanup operations:
\begin{itemize}

\item{} \strong{Plot}: Closes all open graphics devices (except the null device)
\item{} \strong{Global environment}: Removes all objects from the global environment
\item{} \strong{Console}: Clears the console screen (only in interactive sessions)

\end{itemize}


\code{clear()} will not modify the global environment unless you have
previously typed "yes" when prompted. If you do not type "yes", you are asked
again next time; only "yes" is remembered for future sessions.

\strong{Warning}: This function deletes all objects in the global environment.
Save anything that you wish to keep before running.
\end{Details}
%
\begin{Value}
Invisibly returns NULL. Prints a colored confirmation message.
\end{Value}
%
\begin{Examples}
\begin{ExampleCode}

# Interactive use: clear workspace, console, and plots
# First run may prompt; once you type "yes", your preference is saved.
clear()


\end{ExampleCode}
\end{Examples}
\HeaderA{convert\_to\_sql}{Convert CSV file to SQL INSERT statements}{convert.Rul.to.Rul.sql}
%
\begin{Description}
Reads a CSV file and generates SQL statements to insert all rows. Optionally
can also generate a CREATE TABLE statement. The function automatically infers
column types (REAL for numeric, DATE for date strings matching YYYY-MM-DD
format, TEXT otherwise).
\end{Description}
%
\begin{Usage}
\begin{verbatim}
convert_to_sql(input, output, create_table = FALSE)
\end{verbatim}
\end{Usage}
%
\begin{Arguments}
\begin{ldescription}
\item[\code{input}] Character string. Path to the input CSV file.

\item[\code{output}] Character string. Path to the output SQL file where the
statements will be written.

\item[\code{create\_table}] Logical. If \code{TRUE}, includes a CREATE TABLE
statement before the INSERT statements. Default is \code{FALSE}.
\end{ldescription}
\end{Arguments}
%
\begin{Details}
The function performs the following steps:
\begin{enumerate}

\item{} Reads the CSV file using \code{read.csv()} with
\code{stringsAsFactors = FALSE}
\item{} Infers SQL column types:
\begin{itemize}

\item{} Numeric columns become \code{REAL}
\item{} Date columns (matching YYYY-MM-DD format) become \code{DATE}
\item{} All other columns become \code{TEXT}

\end{itemize}

\item{} If \code{create\_table = TRUE}, generates a \code{CREATE TABLE}
statement using the base filename (without extension) as the table name
\item{} Generates \code{INSERT INTO} statements for each row
\item{} Writes all SQL statements to the output file

\end{enumerate}


Single quotes in text values are escaped by doubling them (SQL standard).
Numeric values are inserted without quotes, while text and date values are
wrapped in single quotes.
\end{Details}
%
\begin{Value}
Invisibly returns NULL. The function writes SQL statements to the
specified output file.
\end{Value}
%
\begin{Examples}
\begin{ExampleCode}
# Convert a CSV file to SQL (INSERT statements only)
tmp_csv <- tempfile(fileext = ".csv")
tmp_sql <- tempfile(fileext = ".sql")
write.csv(
  data.frame(id = 1:2, value = c("a", "b"), date = c("2024-01-01", "2024-02-02")),
  tmp_csv,
  row.names = FALSE
)
convert_to_sql(tmp_csv, tmp_sql)

# Convert a CSV file to SQL with CREATE TABLE statement
convert_to_sql(tmp_csv, tmp_sql, create_table = TRUE)

\end{ExampleCode}
\end{Examples}
\HeaderA{desc\_var}{Describe a variable, optionally by groups}{desc.Rul.var}
%
\begin{Description}
Returns a dataframe with one row per group.
\end{Description}
%
\begin{Usage}
\begin{verbatim}
desc_var(y, group = NULL, data = NULL, digits = 3)
\end{verbatim}
\end{Usage}
%
\begin{Arguments}
\begin{ldescription}
\item[\code{y}] A numeric vector of values, a column name (character string or unquoted) if \code{data} is provided,
or a formula of the form \code{y \textasciitilde{} x} or \code{y \textasciitilde{} x1 + x2} (for multiple grouping variables).

\item[\code{group}] Optional grouping variable, if not provided computed for the full data.
Ignored if \code{y} is a formula.

\item[\code{data}] Optional data frame containing the variable(s).

\item[\code{digits}] Number of decimal places to round to. Default is 3.
\end{ldescription}
\end{Arguments}
%
\begin{Details}
The dependent variable (`y`) must be numeric. If you pass an existing
non-numeric variable (e.g., character, factor), `desc\_var()` will stop with
a clear error indicating that the variable must be numeric.
\end{Details}
%
\begin{Value}
A data frame with one row per group (or one row if no group is specified) containing:
\begin{itemize}

\item{} \code{group}: Group identifier
\item{} \code{mean}: Mean
\item{} \code{sd}: Standard deviation
\item{} \code{se}: Standard error
\item{} \code{median}: Median
\item{} \code{min}: Minimum
\item{} \code{max}: Maximum
\item{} \code{mode}: Most frequent value
\item{} \code{freq\_mode}: Frequency of mode
\item{} \code{mode2}: 2nd most frequent value
\item{} \code{freq\_mode2}: Frequency of 2nd mode
\item{} \code{n.total}: Number of observations
\item{} \code{n.missing}: Number of observations with missing (NA) values
\item{} \code{n.unique}: Number of unique values

\end{itemize}

Columns may also carry a \code{"label"} attribute, which provides a short
human-readable description of each column.
\end{Value}
%
\begin{Examples}
\begin{ExampleCode}
# With grouping
df <- data.frame(y = rnorm(100), group = rep(c("A", "B"), 50))
desc_var(y, group, data = df)

# Without grouping (full dataset)
desc_var(y, data = df)

# Direct vectors
y <- rnorm(100)
group <- rep(c("A", "B"), 50)
desc_var(y, group)

# With custom decimal places
desc_var(y, group, data = df, digits = 2)

# Using formula syntax: y ~ x
desc_var(y ~ group, data = df)

# Using formula syntax with multiple grouping variables: y ~ x1 + x2
df2 <- data.frame(y = rnorm(200), x1 = rep(c("A", "B"), 100), x2 = rep(c("X", "Y"), each = 100))
desc_var(y ~ x1 + x2, data = df2)

\end{ExampleCode}
\end{Examples}
\HeaderA{evaluate\_variable\_arguments}{Evaluate Variable Arguments with NSE Support}{evaluate.Rul.variable.Rul.arguments}
\keyword{internal}{evaluate\_variable\_arguments}
%
\begin{Description}
Resolves bare symbols (unquoted names) or quoted strings to actual data,
either from a data frame or the calling environment. This enables both
\code{plot\_density(DV1, data=df)} and \code{plot\_density("DV1", data=df)}
to work identically.
\end{Description}
%
\begin{Usage}
\begin{verbatim}
evaluate_variable_arguments(
  arg_expr,
  arg_name = "arg",
  data = NULL,
  calling_env = parent.frame(),
  func_name = "function",
  allow_null = FALSE
)
\end{verbatim}
\end{Usage}
%
\begin{Arguments}
\begin{ldescription}
\item[\code{arg\_expr}] Unevaluated expression (from match.call() or substitute())

\item[\code{arg\_name}] Character string. Name of the argument (for error messages)

\item[\code{data}] Optional data frame to look up columns

\item[\code{calling\_env}] Environment to evaluate symbols if data is NULL

\item[\code{func\_name}] Character string. Name of calling function (for error messages)

\item[\code{allow\_null}] Logical. If TRUE, NULL is a valid input. Default FALSE.
\end{ldescription}
\end{Arguments}
%
\begin{Value}
A list with:
\begin{itemize}

\item{} \code{value}: The actual data (vector, formula, or NULL)
\item{} \code{name}: Clean name for labels (e.g., "DV1")
\item{} \code{name\_raw}: Raw name for error messages (e.g., "df\$DV1")
\item{} \code{was\_symbol}: Logical. TRUE if input was an unquoted name

\end{itemize}

\end{Value}
\HeaderA{format\_pvalue}{Format P-Values for Display}{format.Rul.pvalue}
%
\begin{Description}
Formats p-values for clean display in figures and tables. e.g., p = .0231, p<.0001
\end{Description}
%
\begin{Usage}
\begin{verbatim}
format_pvalue(p, digits = 4, include_p = FALSE)
\end{verbatim}
\end{Usage}
%
\begin{Arguments}
\begin{ldescription}
\item[\code{p}] A numeric vector of p-values to format.

\item[\code{digits}] Number of decimal places to round to. Default is 4.

\item[\code{include\_p}] Logical. If TRUE, includes "p" prefix before the formatted
value (e.g., "p = .05"). Default is FALSE.
\end{ldescription}
\end{Arguments}
%
\begin{Value}
A character vector of formatted p-values.
\end{Value}
%
\begin{Examples}
\begin{ExampleCode}
# Basic usage
format_pvalue(0.05)
format_pvalue(0.0001)

# More rounding
format_pvalue(0.0001,digits=2)

# Vector input
format_pvalue(c(0.05, 0.001, 0.00001, 0.99))

# With p prefix
format_pvalue(0.05, include_p = TRUE)

\end{ExampleCode}
\end{Examples}
\HeaderA{interprobe}{Probe interactions robustly to nonlinearities}{interprobe}
%
\begin{Description}
Probes an interaction by estimating (or accepting) a model and computing:
- simple slopes ("spotlights") using predicted values
- Johnson-Neyman ("jn") curves using marginal effects
\end{Description}
%
\begin{Usage}
\begin{verbatim}
interprobe(
  x = NULL,
  z = NULL,
  y = NULL,
  model = NULL,
  data = NULL,
  moderator.on.x.axis = TRUE,
  k = NULL,
  spotlights = NULL,
  spotlight.labels = NULL,
  histogram = TRUE,
  max.unique = 11,
  n.bin.continuous = 10,
  n.max = 50,
  xlab = "",
  ylab1 = "",
  ylab2 = "",
  cols = c("red4", "dodgerblue", "green4"),
  main1 = "GAM Simple Slopes",
  main2 = "GAM Johnson-Neyman",
  legend.round = c(1, 4),
  draw = "both",
  save.as = NULL,
  xlim = NULL,
  ylim1 = NULL,
  ylim2 = NULL,
  x.ticks = NULL,
  y1.ticks = NULL,
  y2.ticks = NULL,
  legend.simple.slopes = NULL,
  legend.jn = NULL,
  quiet = FALSE,
  probe.bins = 100
)
\end{verbatim}
\end{Usage}
%
\begin{Arguments}
\begin{ldescription}
\item[\code{x}] The focal predictor. Can be a name (bare or quoted) when `data` or `model`
is provided, or a numeric/factor vector when probing from vectors.

\item[\code{z}] The moderator. Same options as `x`.

\item[\code{y}] The dependent variable. Same options as `x`. Not required when `model` is supplied.

\item[\code{model}] By default `interprobe` estimates a GAM model predicting `y` with `x` and `z`.
You can instead probe a linear interaction by setting model=linear. You can also probe a 
model of your choice by running it separately, saving the output, and submitting it as the model 
argument to  interprobe. This is the way to include covariates for a probed interaction.

\item[\code{data}] Optional data frame containing `x`, `z`, and `y`.

\item[\code{moderator.on.x.axis}] Logical. If TRUE (default), moderator (`z`) is shown on the x-axis.

\item[\code{k}] Integer. Smoothness parameter passed to `mgcv::gam()` when estimating with the default
GAM engine.

\item[\code{spotlights}] Numeric vector of length 3. Values at which curves are computed.

\item[\code{spotlight.labels}] Character vector of length 3. Labels for the legend.

\item[\code{histogram}] Logical. If TRUE (default), show sample size distribution under the plot.

\item[\code{max.unique}] Integer. Threshold for treating a variable as continuous vs discrete.

\item[\code{n.bin.continuous}] Integer. Number of bins used in histogram when binning continuous values.

\item[\code{n.max}] Integer. Sample size at which line darkness/width saturates.

\item[\code{xlab}] Character. X-axis label.

\item[\code{ylab1}] Character. Y-axis label for simple slopes panel.

\item[\code{ylab2}] Character. Y-axis label for JN panel.

\item[\code{cols}] Character vector of length 3. Colors for the three curves.

\item[\code{main1}] Character. Title for simple slopes panel.

\item[\code{main2}] Character. Title for JN panel.

\item[\code{legend.round}] Integer vector length 2. Min/max decimals in legend.

\item[\code{draw}] Which plots to draw: `"both"` (default), `"simple slopes"` (or legacy `"simple.slopes"`), or `"jn"`.

\item[\code{save.as}] Optional file path to save plot (`.png` or `.svg`).

\item[\code{xlim}] Numeric vector length 2. X-axis limits.

\item[\code{ylim1}] Numeric vector length 2. Y-axis limits for simple slopes.

\item[\code{ylim2}] Numeric vector length 2. Y-axis limits for JN.

\item[\code{x.ticks}] Optional custom x-axis ticks.

\item[\code{y1.ticks}] Optional custom y-axis ticks for panel 1.

\item[\code{y2.ticks}] Optional custom y-axis ticks for panel 2.

\item[\code{legend.simple.slopes}] Optional legend title for simple slopes.

\item[\code{legend.jn}] Optional legend title for JN.

\item[\code{quiet}] Logical. If TRUE, reduces console output.

\item[\code{probe.bins}] Integer. Resolution for probing curves (larger = smoother/slower).
\end{ldescription}
\end{Arguments}
%
\begin{Details}
Designed for GAM models but works with any model supported by `marginaleffects`
(including `lm`, `glm`, `mgcv::gam`, and `lm2` / `estimatr::lm\_robust`).
\end{Details}
%
\begin{Value}
Invisibly returns a list with:
\begin{itemize}

\item{} \code{simple.slopes}: data.frame of predicted values and confidence intervals
\item{} \code{johnson.neyman}: data.frame of marginal effects and confidence intervals
\item{} \code{frequencies}: data.frame with bin frequencies used for shading/histogram
\item{} \code{gam\_results}: the fitted GAM model when estimated inside \code{interprobe()}
\item{} \code{gam\_results\_testing}: when \code{interprobe()} estimates a GAM internally and \code{x}
has exactly 2 unique values, a separate GAM fit used for interaction testing (with a \code{ti()}
term and numeric coding of \code{x})
\item{} \code{lm2\_results}: when \code{interprobe()} estimates a GAM internally, also returns the
corresponding linear fit \code{lm2(y \textasciitilde{} x * z)} (or \code{NULL} if package \code{estimatr}
is not installed)

\end{itemize}

\end{Value}
\HeaderA{list2}{Enhanced alternative to list()}{list2}
%
\begin{Description}
List with objects that are automatically named.
\end{Description}
%
\begin{Usage}
\begin{verbatim}
list2(...)
\end{verbatim}
\end{Usage}
%
\begin{Arguments}
\begin{ldescription}
\item[\code{...}] Objects to include in the list. Objects are automatically named 
based on their variable names unless explicit names are provided.
\end{ldescription}
\end{Arguments}
%
\begin{Details}
list2(x , y)      is equivalent to list(x = x , y = y)

list2(x , y2 = y) is equivalent to list(x = x , y2 = y)

Based on: \url{https://stackoverflow.com/questions/16951080/can-lists-be-created-that-name-themselves-based-on-input-object-names}
\end{Details}
%
\begin{Value}
A named list. Each element is named after the variable passed to 
the function (or the explicit name if provided). The structure is identical 
to a standard R list created with \code{\LinkA{list}{list}}.
\end{Value}
%
\begin{Examples}
\begin{ExampleCode}
x <- 1:5
y <- letters[1:3]
z <- matrix(1:4, nrow = 2)

# Create named list from objects
my_list <- list2(x, y, z)
names(my_list)  # "x" "y" "z"

# Works with explicit names too
my_list2 <- list2(a = x, b = y)
names(my_list2)  # "a" "b"

\end{ExampleCode}
\end{Examples}
\HeaderA{lm2}{Enhanced alternative to lm()}{lm2}
%
\begin{Description}
Runs a linear regression with better defaults (robust SE), and richer \& better 
formatted output than \code{lm}. For robust and clustered errors it relies on \code{\LinkA{lm\_robust}{lm.Rul.robust}}. 
The output reports classical and robust errors, number of missing observations per
variable, an effect size column (standardized regression coefficient), and a red.flag column per variable 
flagging the need to conduct specific diagnostics. It relies by default on HC3 for standard errors;
\code{lm\_robust} relies on HC2 (and Stata's 'reg y x, robust' on HC1), which can have
inflated false-positive rates in smaller samples (Long \& Ervin, 2000).
\end{Description}
%
\begin{Arguments}
\begin{ldescription}
\item[\code{se\_type}] The type of standard error to use. Default is \code{"HC3"}.
Without clusters: \code{"HC0"}, \code{"HC1"}, \code{"HC2"}, or \code{"HC3"}.
When \code{clusters} is specified, \code{se\_type} is automatically set to \code{"CR2"}.

\item[\code{notes}] Logical. If TRUE (default), print explanatory notes below the table
when the result is printed.

\item[\code{clusters}] An optional variable indicating clusters for cluster-robust standard 
errors. When specified, \code{se\_type} is automatically set to \code{"CR2"} 
(bias-reduced cluster-robust estimator). Passed to \code{\LinkA{lm\_robust}{lm.Rul.robust}}.

\item[\code{fixed\_effects}] An optional right-sided formula containing the fixed effects 
to be projected out (absorbed) before estimation. Useful for models with many 
fixed effect groups (e.g., \code{\textasciitilde{} firm\_id} or \code{\textasciitilde{} firm\_id + year}). 
Passed to \code{\LinkA{lm\_robust}{lm.Rul.robust}}.

\item[\code{...}] Additional arguments passed to \code{\LinkA{lm\_robust}{lm.Rul.robust}}.
\end{ldescription}
\end{Arguments}
%
\begin{Details}
Robust standard errors and clustered standard errors are computed using 
\code{\LinkA{lm\_robust}{lm.Rul.robust}}; see the documentation of that function for details (using by default CR2 errors)
The output shows both standard errors and when clustering errors it reports all three.
The red.flag column is based on the difference between robust and classical standard errors.

The \code{red.flag} column provides diagnostic warnings:
\begin{itemize}

\item{} \code{!}, \code{!!}, \code{!!!}: Robust and classical standard errors differ by 
more than 25\%, 50\%, or 100\%, respectively. Large differences may suggest model 
misspecification or outliers (but they may also be benign). When encountering a red flag,
authors should plot the distributions to look for outliers or skewed data, and use \code{\LinkA{scatter.gam}{scatter.gam}}
to look for possible nonlinearities in the relevant variables.
King \& Roberts (2015) propose a higher cutoff, at 100\%, and a bootstrapped significance test; 
\code{statuser} does not follow either recommendation. The former seems too liberal, the 
latter too time consuming to include in every regression, plus the focus here is on individual variables rather than joint tests.
\item{} \code{X}: For interaction terms, the component variables are correlated (|r| > 0.3 or p < .05),
which means the interaction term is likely to be biased. See Simonsohn (2024) "Interacting with curves"
\Rhref{https://doi.org/10.1177/25152459231207787}{doi:10.1177\slash{}25152459231207787}.

\end{itemize}

\end{Details}
%
\begin{Value}
An object of class \code{c("lm2", "lm\_robust", "lm")}. This inherits 
from \code{\LinkA{lm\_robust}{lm.Rul.robust}} and can be used with packages like 
\code{marginaleffects}. The object contains all components of an \code{lm\_robust} 
object plus additional attributes:
\begin{description}

\item[statuser\_table] A data frame with columns: \code{term}, \code{estimate}, 
\code{SE.robust}, \code{SE.classical}, \code{t}, \code{df}, \code{p.value}, 
\code{B} (standardized coefficient), and optionally \code{SE.cluster} when 
clustered standard errors are used.
\item[classical\_fit] The underlying \code{lm} object with classical standard errors.
\item[na\_counts] Integer vector of missing value counts per variable.
\item[n\_missing] Total number of observations excluded due to missing values.
\item[has\_clusters] Logical indicating whether clustered standard errors were used.

\end{description}

When printed, displays a formatted regression table with robust and classical 
standard errors, effect sizes, and diagnostic red flags.
\end{Value}
%
\begin{References}
King, G., \& Roberts, M. E. (2015). How robust standard errors expose methodological 
problems they do not fix, and what to do about it. \emph{Political Analysis}, 23(2), 159-179.

Long, J. S., \& Ervin, L. H. (2000). Using heteroscedasticity consistent standard errors 
in the linear regression model. \emph{The American Statistician}, 54(3), 217-224.

Simonsohn, U. (2024). Interacting with curves: How to validly test and probe 
interactions in the real (nonlinear) world. \emph{Advances in Methods and Practices in 
Psychological Science}, 7(1), 1-22.
\end{References}
%
\begin{SeeAlso}
\code{\LinkA{lm\_robust}{lm.Rul.robust}}, \code{\LinkA{scatter.gam}{scatter.gam}}
\end{SeeAlso}
%
\begin{Examples}
\begin{ExampleCode}
# Basic usage with data argument
lm2(mpg ~ wt + hp, data = mtcars)

# Without data argument (variables from environment)
y <- mtcars$mpg
x1 <- mtcars$wt
x2 <- mtcars$hp
lm2(y ~ x1 + x2)

# RED FLAG EXAMPLES

# Example 1: red flag catches a nonlinearity
# True model is quadratic: y = x^2
set.seed(123)
x <- runif(200, -3, 3)
y <- x^2 + rnorm(200, sd = 2)

# lm2() shows red flag due to misspecification
lm2(y ~ x)

# Follow up with scatter.gam() to diagnose it
scatter.gam(x, y)

# Example 2: red flag catches an outlier in y
# True model is y = x, but one observation has a very large y value
set.seed(123)
x <- sort(rnorm(200))
y <- round(x + rnorm(200, sd = 2), 1)
y[200] <- 100  # Outlier

# lm2() flags x
lm2(y ~ x)

# Look at distribution of y to spot the outlier
plot_freq(y)

# Example 3: red flag catches an outlier in one predictor
# True model is y = x1 + x2, but x2 has an extreme value
set.seed(123)
x1 <- round(rnorm(200),.1)
x2 <- round(rnorm(200),.1)
y <- x1 + x2 + rnorm(200, sd = 0.5)
x2[200] <- 50  # Outlier in x2

# lm2() flags x2 (but not x1)
lm2(y ~ x1 + x2)

# Look at distribution of x2 to spot the outlier
plot_freq(x2)

# CLUSTERED STANDARD ERRORS
# When observations are grouped (e.g., students within schools),
# use clusters to account for within-group correlation
set.seed(123)
n_clusters <- 20
n_per_cluster <- 15
cluster_id <- rep(1:n_clusters, each = n_per_cluster)
cluster_effect <- rnorm(n_clusters, sd = 2)[cluster_id]
x <- rnorm(n_clusters * n_per_cluster)
y <- 1 + 0.5 * x + cluster_effect + rnorm(n_clusters * n_per_cluster)
mydata <- data.frame(y = y, x = x, cluster_id = cluster_id)

# Clustered SE (CR2) - note the SE.cluster column in output
lm2(y ~ x, data = mydata, clusters = cluster_id)

# FIXED EFFECTS
# Use fixed_effects to absorb group-level variation (e.g., firm or year effects)
# This is useful for panel data or when you have many fixed effect levels
set.seed(456)
n_firms <- 30
n_years <- 5
firm_id <- rep(1:n_firms, each = n_years)
year <- rep(2018:2022, times = n_firms)
firm_effect <- rnorm(n_firms, sd = 3)[firm_id]
x <- rnorm(n_firms * n_years)
y <- 2 + 0.8 * x + firm_effect + rnorm(n_firms * n_years)
panel <- data.frame(y = y, x = x, firm_id = factor(firm_id), year = factor(year))

# Absorb firm fixed effects (coefficient on x is estimated, firm dummies are not shown)
lm2(y ~ x, data = panel, fixed_effects = ~ firm_id)

# Two-way fixed effects (firm and year)
lm2(y ~ x, data = panel, fixed_effects = ~ firm_id + year)

\end{ExampleCode}
\end{Examples}
\HeaderA{message2}{Enhanced alternative to message()}{message2}
%
\begin{Description}
Add options to set color and to end execution of code (to be used as error message)
\end{Description}
%
\begin{Usage}
\begin{verbatim}
message2(..., col = "cyan", font = 1, stop = FALSE)
\end{verbatim}
\end{Usage}
%
\begin{Arguments}
\begin{ldescription}
\item[\code{...}] Message content to be printed. Multiple arguments are pasted together.

\item[\code{col}] text color. Default is "cyan".

\item[\code{font}] Integer. 1 for plain text (default), 2 for bold text.

\item[\code{stop}] Logical. If TRUE, stops execution (like \code{stop()}) but without printing "Error:".
\end{ldescription}
\end{Arguments}
%
\begin{Details}
This function prints colored messages to the console. If ANSI color codes are supported
by the terminal, the message will be colored. Otherwise, it will be printed as plain text.
If \code{stop = TRUE}, execution will be halted after printing the message.
\end{Details}
%
\begin{Value}
No return value, called for side effects. Prints a colored message 
to the console. If \code{stop = TRUE}, execution is halted after printing 
the message.
\end{Value}
%
\begin{Examples}
\begin{ExampleCode}
message2("This is a plain cyan message", col = "cyan", font = 1)
message2("This is a bold cyan message", col = "cyan", font = 2)
message2("This is a bold red message", col = "red", font = 2)

cat("this will be shown")
try(message2("This stops execution", stop = TRUE), silent = TRUE)
cat("this will be shown after the try")


\end{ExampleCode}
\end{Examples}
\HeaderA{plot\_cdf}{Plot Empirical Cumulative Distribution Functions by Group}{plot.Rul.cdf}
%
\begin{Description}
Plots empirical cumulative distribution functions (ECDFs) separately for
each unique value of a grouping variable, with support for vectorized
plotting parameters. If no grouping variable is provided, plots a single ECDF.
\end{Description}
%
\begin{Usage}
\begin{verbatim}
plot_cdf(
  formula,
  y2 = NULL,
  data = NULL,
  order = NULL,
  show.ks = TRUE,
  show.quantiles = TRUE,
  ...
)
\end{verbatim}
\end{Usage}
%
\begin{Arguments}
\begin{ldescription}
\item[\code{formula}] Two possible uses (similar to \code{t.test()}):
\begin{itemize}

\item{} Single Variable (possibly by subgroup): \code{plot\_cdf(y)} or \code{plot\_cdf(y\textasciitilde{}x)}
\item{} Contrast Two Variables: \code{plot\_cdf(y1, y2)}

\end{itemize}


\item[\code{y2}] optional second variable when contrasting two variables \code{plot\_cdf(y1,y2)}

\item[\code{data}] An optional data frame containing the variables in the formula.
If \code{data} is not provided, variables are evaluated from the calling environment.

\item[\code{order}] Controls the order in which groups appear in the plot and legend. 
Use \code{-1} to reverse the default order. Alternatively, provide a vector specifying
the exact order (e.g., \code{c("B", "A", "C")}). If \code{NULL} (default), groups are 
ordered by their factor levels (if the grouping variable is a factor) or sorted 
alphabetically/numerically. Only applies when using grouped plots.

\item[\code{show.ks}] Logical. If TRUE (default), shows Kolmogorov-Smirnov test results
when there are exactly 2 groups. If FALSE, KS test results are not displayed.

\item[\code{show.quantiles}] Logical. If TRUE (default), shows horizontal lines and results
at 25th, 50th, and 75th percentiles when there are exactly 2 groups. If FALSE,
quantile lines and results are not displayed.

\item[\code{...}] Additional arguments passed to plotting functions. Can be single values
(applied to all groups) or vectors (applied element-wise to each group).
Common parameters include \code{col}, \code{lwd}, \code{lty}, \code{pch},
\code{type}, etc.
\end{ldescription}
\end{Arguments}
%
\begin{Value}
Invisibly returns a list containing:
\begin{itemize}

\item{} \code{ecdfs}: A list of ECDF function objects, one per group. Each can be
called as a function to compute cumulative probabilities (e.g., \code{result\$ecdfs[[1]](5)}
returns P(X <= 5) for group 1).
\item{} \code{ks\_test}: (Only when exactly 2 groups) The Kolmogorov-Smirnov test result
comparing the two distributions. Access p-value with \code{result\$ks\_test\$p.value}.
\item{} \code{quantile\_regression\_25}: (Only when exactly 2 groups) Quantile regression
model for the 25th percentile.
\item{} \code{quantile\_regression\_50}: (Only when exactly 2 groups) Quantile regression
model for the 50th percentile (median).
\item{} \code{quantile\_regression\_75}: (Only when exactly 2 groups) Quantile regression
model for the 75th percentile.
\item{} \code{warnings}: Any warnings captured during execution (if any).

\end{itemize}

\end{Value}
%
\begin{Examples}
\begin{ExampleCode}
# Basic usage with single variable (no grouping)
y <- rnorm(100)
plot_cdf(y)

# Basic usage with formula syntax and grouping
group <- rep(c("A", "B", "C"), c(30, 40, 30))
plot_cdf(y ~ group)

# With custom colors (scalar - same for all)
plot_cdf(y ~ group, col = "blue")

# With custom colors (vector - different for each group)
plot_cdf(y ~ group, col = c("red", "green", "blue"))

# Multiple parameters
plot_cdf(y ~ group, col = c("red", "green", "blue"), lwd = c(1, 2, 3))

# With line type and point character
plot_cdf(y ~ group, col = c("red", "green", "blue"), lty = c(1, 2, 3), lwd = 2)

# Using data frame
df <- data.frame(value = rnorm(100), group = rep(c("A", "B"), 50))
plot_cdf(value ~ group, data = df)
plot_cdf(value ~ group, data = df, col = c("red", "blue"))

# Compare two vectors
y1 <- rnorm(50)
y2 <- rnorm(50, mean = 1)
plot_cdf(y1, y2)

# Formula syntax without data (variables evaluated from environment)
widgetness <- rnorm(100)
gender <- rep(c("M", "F"), 50)
plot_cdf(widgetness ~ gender)

# Using the returned object
df <- data.frame(value = c(rnorm(50, 0), rnorm(50, 1)), group = rep(c("A", "B"), each = 50))
result <- plot_cdf(value ~ group, data = df)

# Use ECDF to find P(X <= 0.5) for group A
result$ecdfs[[1]](0.5)

# Access KS test p-value
result$ks_test$p.value

# Summarize median quantile regression
summary(result$quantile_regression_50)
\end{ExampleCode}
\end{Examples}
\HeaderA{plot\_density}{Plot density of a variable, optionally by another variable}{plot.Rul.density}
%
\begin{Description}
Plots the distribution of a variable by group, simply: \code{plot\_density(y \textasciitilde{} x)}
\end{Description}
%
\begin{Usage}
\begin{verbatim}
plot_density(
  formula,
  y2 = NULL,
  data = NULL,
  order = NULL,
  show_means = TRUE,
  ...
)
\end{verbatim}
\end{Usage}
%
\begin{Arguments}
\begin{ldescription}
\item[\code{formula}] Two possible uses (similar to \code{t.test()}):
\begin{itemize}

\item{} Single Variable (possibly by subgroup): \code{plot\_density(y)} or \code{plot\_density(y\textasciitilde{}x)}
\item{} Contrast Two Variables: \code{plot\_density(y1, y2)}

\end{itemize}


\item[\code{y2}] optional second variable when contrasting two variables \code{plot\_density(y1,y2)}

\item[\code{data}] An optional data frame containing the variables in the formula.

\item[\code{order}] Controls the order in which groups appear in the plot and legend. 
Use \code{-1} to reverse the default order. Alternatively, provide a vector specifying
the exact order (e.g., \code{c("B", "A", "C")}). If \code{NULL} (default), groups are 
ordered by their factor levels (if the grouping variable is a factor) or sorted 
alphabetically/numerically. Only applies when using grouped plots.

\item[\code{show\_means}] Logical. If TRUE (default), shows points at means.

\item[\code{...}] Additional arguments passed to plotting functions.
\end{ldescription}
\end{Arguments}
%
\begin{Details}
Plot parameters like
\code{col}, \code{lwd}, \code{lty}, and \code{pch} can be specified as:
\begin{itemize}

\item{} A single value: applied to all groups
\item{} A vector: applied to groups in order of unique \code{group} values

\end{itemize}

\end{Details}
%
\begin{Value}
Invisibly returns a list with the following element:
\begin{description}

\item[densities] A named list of density objects (class \code{"density"}), 
one for each group. Each density object contains \code{x} (evaluation points), 
\code{y} (density estimates), \code{bw} (bandwidth), and other components 
as returned by \code{\LinkA{density}{density}}. If no grouping variable is 
provided, the list contains a single element named \code{"all"}.

\end{description}

The function is primarily called for its side effect of creating a plot.
\end{Value}
%
\begin{Examples}
\begin{ExampleCode}
# Basic usage with formula syntax (no grouping)
y <- rnorm(100)
plot_density(y)

# With grouping variable
group <- rep(c("A", "B", "C"), c(30, 40, 30))
plot_density(y ~ group)

# With custom colors (scalar - same for all)
plot_density(y ~ group, col = "blue")

# With custom colors (vector - different for each group)
plot_density(y ~ group, col = c("red", "green", "blue"))

# Multiple parameters
plot_density(y ~ group, col = c("red", "green", "blue"), lwd = c(1, 2, 3))

# With line type
plot_density(y ~ group, col = c("red", "green", "blue"), lty = c(1, 2, 3), lwd = 2)

# Using data frame
df <- data.frame(value = rnorm(100), group = rep(c("A", "B"), 50))
plot_density(value ~ group, data = df)
plot_density(value ~ group, data = df, col = c("red", "blue"))

# Compare two vectors
y1 <- rnorm(50)
y2 <- rnorm(50, mean = 1)
plot_density(y1, y2)

\end{ExampleCode}
\end{Examples}
\HeaderA{plot\_freq}{Plot frequencies for a variable (histogram without binning)}{plot.Rul.freq}
%
\begin{Description}
Creates a frequency plot showing the frequency of every observed value, optionaly by group. 
Most frequent values are labeled by default.
\end{Description}
%
\begin{Usage}
\begin{verbatim}
plot_freq(
  formula,
  y2 = NULL,
  data = NULL,
  freq = TRUE,
  order = NULL,
  col = "dodgerblue",
  lwd = 9,
  width = NULL,
  value.labels = "auto",
  ticks.max = 30,
  show.x.value = "auto",
  show.legend = TRUE,
  legend.title = NULL,
  col.text = NULL,
  ...
)
\end{verbatim}
\end{Usage}
%
\begin{Arguments}
\begin{ldescription}
\item[\code{formula}] Two possible uses (similar to \code{t.test()}):
\begin{itemize}

\item{} Single Variable (possibly by subgroup): \code{plot\_freq(y)} or \code{plot\_freq(y\textasciitilde{}x)}
\item{} Contrast Two Variables: \code{plot\_freq(y1, y2)}

\end{itemize}


\item[\code{y2}] optional second variable when contrasting two variables \code{plot\_freq(y1,y2)}

\item[\code{data}] An optional data frame containing the variables in the formula.

\item[\code{freq}] Logical. If TRUE (default), displays frequencies. If FALSE, displays percentages.

\item[\code{order}] Controls the order in which groups appear in the plot and legend. 
Use \code{-1} to reverse the default order. Alternatively, provide a vector specifying
the exact order (e.g., \code{c("B", "A", "C")}). If \code{NULL} (default), groups are 
ordered by their factor levels (if the grouping variable is a factor) or sorted 
alphabetically/numerically. Only applies when using grouped plots or comparing two variables.

\item[\code{col}] Color for the bars.

\item[\code{lwd}] Line width for the frequency bars. Default is 9.

\item[\code{width}] Numeric. Width of the frequency bars. If NULL (default), width is automatically calculated based on the spacing between values.

\item[\code{value.labels}] Controls value labeling. If numeric, shows labels for the \code{value.labels}
highest-frequency values (including ties). Use \code{-1} or \code{"all"} to show all labels and
\code{0} to show none. Use \code{"auto"} to label all values when there are 30 or fewer
unique values; otherwise label the single most frequent value. For backward compatibility, \code{TRUE}
is treated as \code{-1} and \code{FALSE} as \code{0}.

\item[\code{ticks.max}] Integer. Maximum number of unique x values to label on the x-axis. If there are more
than \code{ticks.max} unique values, \code{pretty()} ticks are used instead of labeling every value.

\item[\code{show.x.value}] Either \code{"auto"} (default), \code{TRUE}, or \code{FALSE}. If enabled, draws a
small \code{x=<value>} label just above each frequency value label. When \code{"auto"}, x labels are
shown only when plot\_freq is not already labeling all x values on the x-axis.

\item[\code{show.legend}] Logical. If TRUE (default), displays a legend when \code{group} is specified. If FALSE, no legend is shown.

\item[\code{legend.title}] Character string. Title for the legend when \code{group} is specified. If NULL (default), no title is shown.

\item[\code{col.text}] Color for the value labels. If not specified, uses \code{col} for non-grouped plots or group colors for grouped plots.

\item[\code{...}] Pass on any argument accepted by \code{plot()} e.g., \code{xlab='x-axis'} , \code{main='Distribution of X'}
\end{ldescription}
\end{Arguments}
%
\begin{Details}
This function creates a frequency plot where each observed value is shown
with its frequency. Unlike a standard histogram, there is no binning, unlike
a barplot, non-observed values of the variable are shown with 0 frequency 
instead of skipped.
\end{Details}
%
\begin{Value}
Invisibly returns a data frame with values and their frequencies.
\end{Value}
%
\begin{Examples}
\begin{ExampleCode}
# Simple example
x <- c(1, 1, 2, 2, 2, 5, 5)
plot_freq(x)

# Pass on some common \code{plot()} arguments
plot_freq(x, col = "steelblue", xlab = "Value", ylab = "Frequency",ylim=c(0,7))

# Add to an existing plot
plot_freq(x, col = "dodgerblue")


# Compare two vectors
y1 <- c(1, 1, 2, 2, 2, 5, 5)
y2 <- c(1, 2, 2, 3, 3, 3)
plot_freq(y1, y2)

# Using a data frame with grouping
df <- data.frame(value = c(1, 1, 2, 2, 2, 5, 5), group = c("A", "A", "A", "B", "B", "A", "B"))
plot_freq(value ~ 1, data = df)  # single variable
plot_freq(value ~ group, data = df)  # with grouping

# Control group order in legend and plot
plot_freq(value ~ group, data = df, order = c("B", "A"))  # B first, then A
plot_freq(value ~ group, data = df, order = -1)  # Reverse default order

\end{ExampleCode}
\end{Examples}
\HeaderA{plot\_gam}{Plot GAM Model}{plot.Rul.gam}
%
\begin{Description}
Plots fitted GAM values for focal predictor,
keeping any other predictors in the model at a specified quantile  (default: median)
\end{Description}
%
\begin{Usage}
\begin{verbatim}
plot_gam(
  model,
  predictor,
  quantile.others = 50,
  col = "blue4",
  bg = adjustcolor("dodgerblue", 0.2),
  plot2 = "auto",
  col2 = NULL,
  bg2 = "gray90",
  ...
)
\end{verbatim}
\end{Usage}
%
\begin{Arguments}
\begin{ldescription}
\item[\code{model}] A GAM model object fitted using \code{mgcv::gam()}.

\item[\code{predictor}] Character string specifying the name of the predictor variable
to plot on the x-axis.

\item[\code{quantile.others}] Number between 1 and 99 for quantile
at which all other predictors are held constant. Default is 50 (median).

\item[\code{col}] Color for the prediction line. Default is "blue4".

\item[\code{bg}] Background color for the confidence band. Default is
\code{adjustcolor('dodgerblue', .2)}.

\item[\code{plot2}] How to plot the distribution in the lower plot. Options: \code{'auto'} (default,
auto-select based on number of unique values), \code{'freq'} (always plot frequencies),
\code{'density'} (always plot the density) or \code{'none'} (neither). When \code{'auto'}, 
plots frequencies with predictor has less than 30 unique values, density otherwise.

\item[\code{col2}] Color for the lines/bars in the bottom distribution plot. Default is "dodgerblue"

\item[\code{bg2}] Background color for the bottom distribution plot. Default is "gray90".

\item[\code{...}] Additional arguments passed to \code{plot()} and \code{lines()}.
\end{ldescription}
\end{Arguments}
%
\begin{Value}
Invisibly returns a list containing:
\begin{itemize}

\item{} \code{predictor\_values}: The sequence of predictor values used
\item{} \code{predicted}: The predicted values
\item{} \code{se}: The standard errors
\item{} \code{lower}: Lower confidence bound (predicted - 2*se)
\item{} \code{upper}: Upper confidence bound (predicted + 2*se)

\end{itemize}

\end{Value}
%
\begin{Examples}
\begin{ExampleCode}

library(mgcv)
# Fit a GAM model
data(mtcars)
mtcars$cyl <- factor(mtcars$cyl)  # Convert to factor before fitting GAM
model <- gam(mpg ~ s(hp) + s(wt) + cyl, data = mtcars)

# Plot effect of hp (with other variables at median)
plot_gam(model, "hp")

# Plot effect of hp (with other variables at 25th percentile)
plot_gam(model, "hp", quantile.others = 25)

# Customize plot
plot_gam(model, "hp", main = "Effect of Horsepower", col = "blue", lwd = 2)


\end{ExampleCode}
\end{Examples}
\HeaderA{plot\_means}{Barplot of means}{plot.Rul.means}
%
\begin{Description}
Plots means, with confidence intervals, and (optionally) p-values for differences of means and interactions
\end{Description}
%
\begin{Usage}
\begin{verbatim}
plot_means(
  formula,
  data = NULL,
  cluster = NULL,
  tests = "auto",
  quiet = FALSE,
  order = NULL,
  legend.title = NULL,
  col = NULL,
  col.text = NULL,
  values.cex = 1,
  values.align = "top",
  values.round = 1,
  pvalue.cex = 0.9,
  pvalue.col = "gray50",
  ci.level = 95,
  buffer.top = "auto",
  ...
)
\end{verbatim}
\end{Usage}
%
\begin{Arguments}
\begin{ldescription}
\item[\code{formula}] A formula, e.g., \code{y \textasciitilde{} x1+x2}, where x1 \& x2 are
grouping variables (e.g., condition indicators in a 2x2 experiment). 
The formula can include up to three grouping variables. The plot will 
show contiguous bars for x1, and sets of bars for x2 and x3

\item[\code{data}] Optional data frame containing variables in the formula.

\item[\code{cluster}] Optional clustering variable when there are repeated 
observations per cluster 
e.g., \code{cluster="participant\_ID"}. 
When provided, inference for reported tests is based on regressions with
clustered standard errors (via \code{lm2(..., clusters=...)}).

\item[\code{tests}] specifies which comparisons of means to report. Syntax involves 
putting column numbers in a character string, with a - to symbolize a 
comparison and a + symbolizing combination. For example tests="1-2" reports 
t-test comparing columns 1 and 2. tests="1-2,3-4" t-tests comparing columns 
1 \& 2 and another 3 \& 4. To run an interaction use parentheses:
tests=\code{"(4-3)-(2-1)"} and can also be combined with simple tests. 
Main effects can be specified using `+`, for example (1+2)-(3+4) compares 
all observations in the first two columns with all observations in 
the next two columns.

\item[\code{quiet}] Logical. When \code{TRUE}, suppresses console messages from \code{plot\_means()}.

\item[\code{order}] Controls the order of \code{x1} groups (bar order and colors).
Use \code{-1} to reverse the default order (e.g., if plot shows 'male' first and 'female' second, order=-1 will flip that).

\item[\code{legend.title}] Character string. Title for the legend. If \code{NULL},
no title is shown.

\item[\code{col}] Color(s) for \code{x1} bars. If \code{NULL}, colors are chosen
automatically.

\item[\code{col.text}] Color for confidence intervals and other non-bar annotations.
If \code{NULL}, defaults to a dark gray.

\item[\code{values.cex}] Numeric scalar controlling text size for mean value labels
(and related annotations).

\item[\code{values.align}] Where within the bars to put the mean value labels: \code{"top"}, \code{"middle"},
\code{"bottom"}, or \code{"none"}.

\item[\code{values.round}] Non-negative integer. Number of decimal places for mean
value labels.

\item[\code{pvalue.cex}] Numeric scalar controlling p-value label size.

\item[\code{pvalue.col}] Color for p-value brackets/labels.

\item[\code{ci.level}] Confidence interval level for \code{ciL}/\code{ciH} in \code{\$means}.
Default is \code{95}. You can also pass a proportion (e.g., \code{0.95}).

\item[\code{buffer.top}] Either \code{"auto"} (default) or a numeric value. Extra
vertical headroom (as a fraction of the data y-range) added above the
maximum y value to make room for annotations. When \code{"auto"}, uses 0.35
when an interaction p-value is shown (scenario 2) and 0.25 otherwise.

\item[\code{...}] Additional arguments passed to \code{plot()} (e.g., \code{main},
\code{ylim}, \code{ylab}).
\end{ldescription}
\end{Arguments}
%
\begin{Details}
When \code{tests="auto"}, the function reports a small default set of
differences-of-means tests (when applicable) and, in 2x2 designs, an
interaction test:
\begin{description}

\item[Differences in means] If \code{cluster} is \code{NULL}, these are Welch
two-sample t-tests computed with \code{t.test(..., var.equal=FALSE)}. If
\code{cluster} is provided, these comparisons are computed from a
regression using \code{lm2()} with clustered standard errors.
\item[Interaction] The interaction is tested using a linear regression fit
with \code{lm2()}, even when \code{cluster} is \code{NULL}; when
\code{cluster} is provided, the interaction test uses clustered standard
errors.

\end{description}

The regression-based tests use heteroskedasticity-robust inference (HC3) when
\code{cluster} is \code{NULL}. HC3 is a common small-sample adjustment to
White-type robust standard errors and is used to reduce sensitivity to
heteroskedasticity. When \code{cluster} is provided, \code{plot\_means()}
instead uses clustered standard errors (robust to within-cluster correlation).

In the returned \code{\$means} table, \code{ciL} and \code{ciH} are the lower and
upper bounds of a \code{ci.level}\% confidence interval for the mean (when
available). The same confidence level is used for the confidence-interval
whiskers drawn in the figure.
\end{Details}
%
\begin{Value}
A minimal list returned invisibly with two elements:
\begin{description}

\item[\code{means}] A data frame of means (and, when available, confidence intervals)
aligned to the plotting grid.
\item[\code{tests}] A data frame of comparisons used for p-value
annotation (or \code{NULL} if not applicable).

\end{description}

\end{Value}
%
\begin{Examples}
\begin{ExampleCode}
df <- data.frame(y = rnorm(100), group = rep(c("A", "B"), 50))
plot_means(y ~ group, data = df)

df2 <- data.frame(
  y = rnorm(200),
  x1 = rep(c("A", "B"), 100),
  x2 = rep(c("X", "Y"), each = 100)
)
plot_means(y ~ x1 + x2, data = df2)

df3 <- data.frame(
  y = rnorm(600),
  x1 = rep(c("control", "treatment"), times = 300),
  x2 = rep(rep(c("low", "high"), each = 150), times = 2),
  x3 = rep(c("online", "lab"), each = 300)
)
plot_means(y ~ x1 + x2 + x3, data = df3)

\end{ExampleCode}
\end{Examples}
\HeaderA{predict.lm2}{Predict method for lm2 objects}{predict.lm2}
%
\begin{Description}
Predict method for lm2 objects
\end{Description}
%
\begin{Usage}
\begin{verbatim}
## S3 method for class 'lm2'
predict(object, newdata, ...)
\end{verbatim}
\end{Usage}
%
\begin{Arguments}
\begin{ldescription}
\item[\code{object}] An object of class \code{lm2}

\item[\code{newdata}] An optional data frame in which to look for variables with 
which to predict. If omitted, the original model data is used.

\item[\code{...}] Additional arguments passed to \code{\LinkA{predict.lm\_robust}{predict.lm.Rul.robust}},
including \code{se.fit} and \code{interval}.
\end{ldescription}
\end{Arguments}
%
\begin{Value}
A vector of predicted values (or a list with \code{fit} and \code{se.fit} 
if \code{se.fit = TRUE}, or a matrix with \code{fit}, \code{lwr}, \code{upr} 
if \code{interval} is specified)
\end{Value}
\HeaderA{print.desc\_var}{Print method for desc\_var objects}{print.desc.Rul.var}
%
\begin{Description}
Print method for desc\_var objects
\end{Description}
%
\begin{Usage}
\begin{verbatim}
## S3 method for class 'desc_var'
print(x, ...)
\end{verbatim}
\end{Usage}
%
\begin{Arguments}
\begin{ldescription}
\item[\code{x}] An object of class \code{desc\_var}

\item[\code{...}] Additional arguments passed to print.data.frame
\end{ldescription}
\end{Arguments}
%
\begin{Value}
Invisibly returns the original object
\end{Value}
\HeaderA{print.lm2}{Print method for lm2 objects}{print.lm2}
%
\begin{Description}
Print method for lm2 objects
\end{Description}
%
\begin{Usage}
\begin{verbatim}
## S3 method for class 'lm2'
print(x, notes = NULL, ...)
\end{verbatim}
\end{Usage}
%
\begin{Arguments}
\begin{ldescription}
\item[\code{x}] An object of class \code{lm2}

\item[\code{notes}] Logical. If TRUE (default), print explanatory notes below the table.
If not specified, uses the value set when \code{lm2()} was called.

\item[\code{...}] Additional arguments (ignored)
\end{ldescription}
\end{Arguments}
%
\begin{Value}
Invisibly returns the original object
\end{Value}
\HeaderA{print.t.test2}{Print method for t.test2 output}{print.t.test2}
%
\begin{Description}
Print method for t.test2 output
\end{Description}
%
\begin{Usage}
\begin{verbatim}
## S3 method for class 't.test2'
print(x, ...)
\end{verbatim}
\end{Usage}
%
\begin{Arguments}
\begin{ldescription}
\item[\code{x}] An object of class \code{t.test2}

\item[\code{...}] Additional arguments passed to print
\end{ldescription}
\end{Arguments}
%
\begin{Value}
Invisibly returns the input object \code{x}. Called for its side effect 
of printing a formatted t-test summary to the console, including means, 
confidence intervals, test statistics, p-values, sample sizes, and 
APA-formatted results.
\end{Value}
\HeaderA{print.table2}{Print method for table2 output with centered column variable name}{print.table2}
%
\begin{Description}
Print method for table2 output with centered column variable name

Print method for table2 objects
\end{Description}
%
\begin{Usage}
\begin{verbatim}
## S3 method for class 'table2'
print(x, ...)

## S3 method for class 'table2'
print(x, ...)
\end{verbatim}
\end{Usage}
%
\begin{Arguments}
\begin{ldescription}
\item[\code{x}] An object of class \code{table2}

\item[\code{...}] Additional arguments (ignored)
\end{ldescription}
\end{Arguments}
%
\begin{Value}
Invisibly returns the input object \code{x}. Called for its side effect 
of printing a formatted cross-tabulation table to the console. The output 
includes frequencies, optional relative frequencies (row, column, or overall 
proportions), and chi-squared test results when applicable.

Invisibly returns the original object
\end{Value}
\HeaderA{reg2}{Two-Line Interrupted Regression}{reg2}
\keyword{internal}{reg2}
%
\begin{Description}
Performs an interrupted regression analysis with heteroskedastic robust standard errors.
This function fits two regression lines with a breakpoint at a specified value of the
first predictor variable.
\end{Description}
%
\begin{Usage}
\begin{verbatim}
reg2(f, xc, graph = 1, family = "gaussian", data = NULL)
\end{verbatim}
\end{Usage}
%
\begin{Arguments}
\begin{ldescription}
\item[\code{f}] A formula object specifying the model (e.g., y \textasciitilde{} x1 + x2 + x3).
The first predictor is the one on which the u-shape is tested.

\item[\code{xc}] Numeric value specifying where to set the breakpoint.

\item[\code{graph}] Integer. If 1 (default), produces a plot. If 0, no plot is generated.

\item[\code{family}] Character string specifying the family for the GLM model.
Default is "gaussian" for OLS regression. Use "binomial" for probit models.

\item[\code{data}] A data frame containing the variables in the formula.
\end{ldescription}
\end{Arguments}
%
\begin{Details}
This function fits two interrupted regression lines with heteroskedastic robust
standard errors using the sandwich package. The first predictor variable is split
at the breakpoint \code{xc}, creating separate slopes before and after the breakpoint.
\end{Details}
%
\begin{Value}
A list containing:
\begin{itemize}

\item{} \code{b1}, \code{b2}: Slopes of the two regression lines
\item{} \code{p1}, \code{p2}: P-values for the slopes
\item{} \code{z1}, \code{z2}: Z-statistics for the slopes
\item{} \code{u.sig}: Indicator (0/1) for whether u-shape is significant
\item{} \code{xc}: The breakpoint value
\item{} \code{glm1}, \code{glm2}: The fitted GLM models
\item{} \code{rob1}, \code{rob2}: Robust coefficient test results
\item{} \code{msg}: Messages about standard error computation
\item{} \code{yhat.smooth}: Fitted smooth values (if graph=1)

\end{itemize}

\end{Value}
\HeaderA{resize\_images}{Resize Images}{resize.Rul.images}
%
\begin{Description}
Saves images to PNG with a specified width. As input it accepts (SVG, PDF, EPS, JPG, JPEG, TIF, TIFF, BMP, PNG)
Saves to subdirectory '/resized' within input folder (or same directory as file if input is a single file)
\end{Description}
%
\begin{Usage}
\begin{verbatim}
resize_images(path, width)
\end{verbatim}
\end{Usage}
%
\begin{Arguments}
\begin{ldescription}
\item[\code{path}] Character string. Path to a folder containing image files, or path to a single image file.

\item[\code{width}] Numeric vector. Target width(s) in pixels for the output PNG files.
Can be a single value (recycled for all files) or a vector matching the number
of files found.
\end{ldescription}
\end{Arguments}
%
\begin{Details}
This function:
\begin{itemize}

\item{} Searches for image files with extensions: svg, pdf, eps, jpg, jpeg, tif, tiff, bmp, png
\item{} Creates a "resized" subfolder in the target directory if it doesn't exist
\item{} Converts each file to PNG format at the specified width(s)
\item{} Saves output files as: \code{originalname\_width.png} in the resized subfolder

\end{itemize}


Supported input formats:
\begin{itemize}

\item{} Vector graphics: SVG, PDF, EPS (rasterized using rsvg/magick)
\item{} Raster images: JPG, JPEG, TIF, TIFF, BMP, PNG

\end{itemize}

\end{Details}
%
\begin{Value}
Invisibly returns TRUE on success.
\end{Value}
%
\begin{Note}
Dependencies required: \code{rsvg}, \code{magick}, and \code{tools} (base R).
SVG files are rasterized using \code{rsvg::rsvg()}, while PDF/EPS and other
formats are handled by \code{magick::image\_read()}.
\end{Note}
%
\begin{Examples}
\begin{ExampleCode}

# Create a temporary PNG file and resize it
tmp_png <- tempfile(fileext = ".png")
grDevices::png(tmp_png, width = 400, height = 300)
old_par <- graphics::par(no.readonly = TRUE)
graphics::par(mar = c(2, 2, 1, 1))
graphics::plot(1:2, 1:2, type = "n")
grDevices::dev.off()
graphics::par(old_par)
resize_images(tmp_png, width = 80)


\end{ExampleCode}
\end{Examples}
\HeaderA{scatter.gam}{Scatter Plot with GAM Smooth Line}{scatter.gam}
%
\begin{Description}
Creates a scatter plot with a GAM (Generalized Additive Model) smooth line.
Supports both \code{scatter.gam(x, y)} and \code{scatter.gam(y \textasciitilde{} x)}.
\end{Description}
%
\begin{Usage}
\begin{verbatim}
scatter.gam(
  x,
  y,
  data.dots = TRUE,
  three.dots = FALSE,
  data = NULL,
  k = NULL,
  plot.dist = NULL,
  dot.pch = 16,
  dot.col = adjustcolor("gray", 0.7),
  jitter = FALSE,
  ...
)
\end{verbatim}
\end{Usage}
%
\begin{Arguments}
\begin{ldescription}
\item[\code{x}] A numeric vector of x values, or a formula of the form \code{y \textasciitilde{} x}.

\item[\code{y}] A numeric vector of y values. Not used if \code{x} is a formula.

\item[\code{data.dots}] Logical. If TRUE, displays data on scatterplot

\item[\code{three.dots}] Logical. If TRUE, divides x into tertiles and puts markers on the average  x \& y for each

\item[\code{data}] An optional data frame containing the variables \code{x} and \code{y}.

\item[\code{k}] Optional integer specifying the basis dimension for the smooth term
in the GAM model (passed to \code{s(x, k=k)}). If NULL (default), uses the
default basis dimension.

\item[\code{plot.dist}] Character string specifying how to plot the distribution of \code{x}
underneath the scatter plot. Options: \code{NULL} (default, auto-select based on
number of unique values), \code{"none"} (no distribution plot), \code{"plot\_freq"}
(always use \code{plot\_freq()}), or \code{"hist"} (always use \code{hist()}).
When \code{NULL}, uses \code{plot\_freq()} if there are 25 or fewer unique values,
otherwise uses \code{hist()}.

\item[\code{dot.pch}] Plotting character for data points when \code{data.dots = TRUE}.
Default is 16 (filled circle).

\item[\code{dot.col}] Color for data points when \code{data.dots = TRUE}. Default is
\code{adjustcolor('gray', 0.7)} (semi-transparent gray).

\item[\code{jitter}] Logical. If TRUE, applies a small amount of jitter to data points
to reduce overplotting. Default is FALSE.

\item[\code{...}] Additional arguments passed to \code{plot()} and \code{gam()}.
Common plot arguments include:
\begin{itemize}

\item{} \code{main}: Custom title for the plot (e.g., \code{main = "My Title"})
\item{} \code{col}: Color of the GAM smooth line (e.g., \code{col = "red"})
\item{} \code{lwd}: Line width of the GAM smooth line (e.g., \code{lwd = 2})
\item{} \code{xlim}, \code{ylim}: Axis limits (e.g., \code{xlim = c(0, 10)})
\item{} \code{xlab}, \code{ylab}: Axis labels (e.g., \code{xlab = "Age"})

\end{itemize}

\end{ldescription}
\end{Arguments}
%
\begin{Details}
This function fits a GAM model with a smooth term for x and plots the fitted
smooth line. The function uses the \code{mgcv} package's \code{gam()} function.

When \code{three.dots = TRUE}, the x variable is divided into three equal-sized
groups (tertiles), and the mean x and y values for each group are plotted as
points. This provides a simple summary of the relationship across the range of x.
\end{Details}
%
\begin{Value}
Invisibly returns the fitted GAM model object.
\end{Value}
%
\begin{SeeAlso}
\code{\LinkA{scatter.smooth}{scatter.smooth}} for a simpler loess-based scatter plot smoother.
\end{SeeAlso}
%
\begin{Examples}
\begin{ExampleCode}
# Generate sample data for examples
x <- rnorm(100)
y <- 2*x + rnorm(100)

# Plot GAM smooth line only
scatter.gam(x, y)

# Equivalent call using formula syntax (y ~ x)
scatter.gam(y ~ x)

# Include scatter plot with underlying data points behind the GAM line
scatter.gam(x, y, data.dots = TRUE)

# Include summary points showing mean x and y for each tertile bin
scatter.gam(x, y, three.dots = TRUE)

# Customize the plot with a custom title, line color, and line width
scatter.gam(x, y, data.dots = TRUE, col = "red", lwd = 2, main = "GAM Fit")

# Control smoothness of the GAM line by specifying the basis dimension
scatter.gam(x, y, k = 10)

\end{ExampleCode}
\end{Examples}
\HeaderA{summary.lm2}{Summary method for lm2 objects}{summary.lm2}
%
\begin{Description}
Summary method for lm2 objects
\end{Description}
%
\begin{Usage}
\begin{verbatim}
## S3 method for class 'lm2'
summary(object, ...)
\end{verbatim}
\end{Usage}
%
\begin{Arguments}
\begin{ldescription}
\item[\code{object}] An object of class \code{lm2}

\item[\code{...}] Additional arguments passed to \code{\LinkA{print.lm2}{print.lm2}}
\end{ldescription}
\end{Arguments}
%
\begin{Value}
Invisibly returns the original object
\end{Value}
\HeaderA{t.test2}{Enhanced alternative to t.test()}{t.test2}
%
\begin{Description}
The basic t-test function in R, \code{\LinkA{t.test}{t.test}}, does not report 
the observed difference of means, does not stipulate which mean
is subtracted from which (i.e., whether it computed A-B or B-A), and
presents the test results on the console in a verbose unorganized 
paragraph of text. \code{t.test2} improves on all those counts, and 
in addition, it reports the number of observations per group and if any observations 
are missing it issues a warning. It returns a dataframe instead of a list.
The function is also registered as \code{t.test2} for the \code{\LinkA{t}{t}}
generic to satisfy S3 registration checks while keeping direct calls to
\code{t.test2(...)} unchanged.
\end{Description}
%
\begin{Arguments}
\begin{ldescription}
\item[\code{x}] The first data argument passed to \code{\LinkA{t.test}{t.test}}. This can
be a numeric vector or a formula.

\item[\code{...}] Arguments passed to \code{\LinkA{t.test}{t.test}}
\end{ldescription}
\end{Arguments}
%
\begin{Value}
A data frame with class \code{c("t.test2", "data.frame")} containing 
a single row with the following columns:
\begin{description}

\item[mean columns] One or two columns containing group means, named after 
the input variables (e.g., \code{men}, \code{women}) or \code{Group 1}, 
\code{Group 2} for long names.
\item[diff column] For two-sample tests, the difference between means 
(e.g., \code{men-women}).
\item[ci] The confidence level as a string (e.g., "95 percent").
\item[ci.L, ci.H] Lower and upper bounds of the confidence interval.
\item[t] The t-statistic.
\item[df] Degrees of freedom.
\item[p.value] The p-value.
\item[N columns] Sample sizes, named \code{N(group1)}, \code{N(group2)} or 
\code{N1}, \code{N2}. For paired tests, a single \code{N} column.
\item[correlation] For paired tests only, the correlation between pairs.

\end{description}

Attributes store additional information including missing value counts and 
test type (one-sample, two-sample, paired, Welch vs. Student).
\end{Value}
%
\begin{Examples}
\begin{ExampleCode}
# Two-sample t-test
men <- rnorm(100, mean = 5, sd = 1)
women <- rnorm(100, mean = 4.8, sd = 1)
t.test2(men, women)

# Paired t-test
x <- rnorm(50, mean = 5, sd = 1)
y <- rnorm(50, mean = 5.2, sd = 1)
t.test2(x, y, paired = TRUE)

# One-sample t-test
data <- rnorm(100, mean = 5, sd = 1)
t.test2(data, mu = 0)

# Formula syntax
data <- data.frame(y = rnorm(100), group = rep(c("A", "B"), 50))
t.test2(y ~ group, data = data)

\end{ExampleCode}
\end{Examples}
\HeaderA{table2}{Enhanced alternative to table()}{table2}
%
\begin{Description}
The function \code{\LinkA{table}{table}} does not show variable 
names when tabulating from a dataframe, requires running another
function, \code{\LinkA{prop.table}{prop.table}},  to tabulate proportions 
and yet another function, \code{\LinkA{chisq.test}{chisq.test}} to test difference of 
proportions.  \code{table2} does what those three functions do, producing easier to 
read output, and always shows variable names.
\end{Description}
%
\begin{Arguments}
\begin{ldescription}
\item[\code{...}] same arguments as \code{\LinkA{table}{table}}, plus the arguments shown below

\item[\code{prop}] report a table with:
\begin{itemize}

\item{} \code{prop="all"}: Proportions for full table (each cell / total)
\item{} \code{prop="row"}: Proportions by row  ('rows' also accepted)
\item{} \code{prop="col"}: Proportions by columns ('cols', 'column', 'columns' also accepted)

\end{itemize}


\item[\code{digits}] Number of decimal values to show for proportions

\item[\code{chi}] Logical. If \code{TRUE}, performs a chi-square test on frequency table,
reports results in APA format

\item[\code{correct}] Logical. If \code{TRUE}, applies Yates' continuity correction 
for 2x2 tables in the chi-square test. Default is \code{FALSE} (no correction).
\end{ldescription}
\end{Arguments}
%
\begin{Value}
A list (object of class "table2") with the following components:
\begin{itemize}

\item{} \code{freq}: frequency table 
\item{} \code{prop}: proportions table 
\item{} \code{chisq}: chi-square test  

\end{itemize}

\end{Value}
%
\begin{Examples}
\begin{ExampleCode}
# Create example data
df <- data.frame(
  group = c("A", "A", "B", "B", "A"),
  status = c("X", "Y", "X", "Y", "X")
)

# Enhanced table with variable names (2 variables)
table2(df$group, df$status)

# Enhanced table with variable names (3 variables)
df3 <- data.frame(
  x = c("A", "A", "B", "B"),
  y = c("X", "Y", "X", "Y"),
  z = c("high", "low", "high", "low")
)
table2(df3$x, df3$y, df3$z)

# Table with proportions
table2(df$group, df$status, prop = 'all')  # Overall proportions
table2(df$group, df$status, prop = 'row')  # Row proportions
table2(df$group, df$status, prop = 'col')  # Column proportions

# Table with chi-square test
table2(df$group, df$status, chi = TRUE,prop='all')

\end{ExampleCode}
\end{Examples}
\HeaderA{text2}{Enhanced alternative to text()}{text2}
%
\begin{Description}
Adds to text() optional background color and verbal alignment (align='center')
\end{Description}
%
\begin{Arguments}
\begin{ldescription}
\item[\code{x}, \code{y}] coordinates for text placement

\item[\code{labels}] text to display

\item[\code{align}] alignment in relation to x coordinate ('left','center','right')

\item[\code{bg}] background color

\item[\code{cex}] character expansion factor

\item[\code{pad}] left/right padding in percentage (e.g., .03)

\item[\code{pad\_v}] top/bottom padding in percentage (e.g., .25)

\item[\code{...}] Additional arguments passed to \code{\LinkA{text}{text}}.
\end{ldescription}
\end{Arguments}
%
\begin{Value}
No return value, called for side effects. Adds text with an optional 
background rectangle to an existing plot.
\end{Value}
%
\begin{Examples}
\begin{ExampleCode}
# Create a simple plot
plot(1:10, 1:10, type = "n", main = "text2() - Alignment & Color")

# Alignment respect to x=5
text2(5, 8, "align='left' from 5", align = "left", bg = "yellow1")
text2(5, 7, "align='right' from 5", align = "right", bg = "blue", col = "white")
text2(5, 6, "align='center' from 5", align = "center", bg = "black", col = "white")
abline(v = 5, lty = 2)

# Multiple labels with different alignments
text2(c(2, 5, 8), c(5, 5, 5), 
      labels = c("Left", "Center", "Right"),
      align = c("left", "center", "right"),
      bg = c("pink", "lightblue", "lightgreen"))

# Text with custom font color (passed through ...)
text2(5, 3, "Red Text", col = "red", bg = "white")

# Padding examples
plot(1:10, 1:10, type = "n", main = "Padding Examples")

# Default padding (pad=0.03, pad_v=0.25)
text2(5, 8, "Default padding", bg = "lightblue")

# More horizontal padding
text2(5, 6, "Wide padding", pad = 0.2, bg = "lightgreen")

# More vertical padding
text2(5, 4, "Tall padding", pad_v = 0.8, bg = "lightyellow")

# Both padding increased
text2(5, 2, "Extra padding", pad = 0.15, pad_v = 0.6, bg = "pink")
\end{ExampleCode}
\end{Examples}
\HeaderA{twolines}{Two-Lines Test of U-Shapes}{twolines}
%
\begin{Description}
Implements the two-lines test for U-shaped (or inverted U-shaped) relationships 
introduced by Simonsohn (2018).
\end{Description}
%
\begin{Usage}
\begin{verbatim}
twolines(
  f,
  graph = 1,
  link = "gaussian",
  data = NULL,
  pngfile = "",
  quiet = FALSE
)
\end{verbatim}
\end{Usage}
%
\begin{Arguments}
\begin{ldescription}
\item[\code{f}] A formula object specifying the model (e.g., y \textasciitilde{} x1 + x2 + x3).
The first predictor is the one tested for a u-shaped relationship.

\item[\code{graph}] Integer. If 1 (default), produces a plot. If 0, no plot is generated.

\item[\code{link}] Character string specifying the link function for the GAM model.
Default is "gaussian".

\item[\code{data}] An optional data frame containing the variables in the formula.
If not provided, variables are evaluated from the calling environment.

\item[\code{pngfile}] Optional character string. If provided, saves the plot to a PNG file
with the specified filename.

\item[\code{quiet}] Logical. If TRUE, suppresses the Robin Hood details messages.
Default is FALSE.
\end{ldescription}
\end{Arguments}
%
\begin{Details}
Reference: 
Simonsohn, Uri (2018) "Two lines: A valid alternative to the invalid testing of 
U-shaped relationships with quadratic regressions." AMPPS, 538-555. 
\Rhref{https://doi.org/10.1177/2515245918805755}{doi:10.1177\slash{}2515245918805755}


The test beings fitting a GAM model, predicting y with a smooth of x, and optionally with covariates. 
It identifies the interior most extreme value of fitted y, and adjusts from the matching x-value 
to set the breakpoint relying on the Robin Hood procedure introduced also by Simonsohn (2018).
It then estimates the (once) interrupted regression using that breakpoint,
and reports the slope and significance of the average slopes at either side of it. A U-shape 
is significant if the slopes are of opposite sign and are both individually significant.
\end{Details}
%
\begin{Value}
A list containing:
\begin{itemize}

\item{} All elements from \code{reg2()}: \code{b1}, \code{b2}, \code{p1}, \code{p2},
\code{z1}, \code{z2}, \code{u.sig}, \code{xc}, \code{glm1}, \code{glm2}, \code{rob1},
\code{rob2}, \code{msg}, \code{yhat.smooth}
\item{} \code{yobs}: Observed y values (adjusted for covariates if present)
\item{} \code{y.hat}: Fitted values from GAM
\item{} \code{y.ub}, \code{y.lb}: Upper and lower bounds for fitted values
\item{} \code{y.most}: Most extreme fitted value
\item{} \code{x.most}: x-value associated with most extreme fitted value
\item{} \code{f}: Formula as character string
\item{} \code{bx1}, \code{bx2}: Linear and quadratic coefficients from preliminary quadratic regression
\item{} \code{minx}: Minimum x value
\item{} \code{midflat}: Median of flat region
\item{} \code{midz1}, \code{midz2}: Z-statistics at midpoint

\end{itemize}

\end{Value}
%
\begin{Examples}
\begin{ExampleCode}

# Simple example with simulated data
set.seed(123)
x <- rnorm(100)
y <- -x^2 + rnorm(100)
data <- data.frame(x = x, y = y)
result <- twolines(y ~ x, data = data)

# With covariates
z <- rnorm(100)
y <- -x^2 + 0.5*z + rnorm(100)
data <- data.frame(x = x, y = y, z = z)
result <- twolines(y ~ x + z, data = data)

# Without data argument (variables evaluated from environment)
x <- rnorm(100)
y <- -x^2 + rnorm(100)
result <- twolines(y ~ x)


\end{ExampleCode}
\end{Examples}
\HeaderA{validate\_formula}{Validate Formula Variables}{validate.Rul.formula}
\keyword{internal}{validate\_formula}
%
\begin{Description}
Checks if the input is a formula and validates that all variables mentioned
in the formula exist either in the provided data frame or in the environment.
This is a lightweight validation function that should be called early in functions
that accept formula syntax.
\end{Description}
%
\begin{Usage}
\begin{verbatim}
validate_formula(
  formula,
  data = NULL,
  func_name = "function",
  calling_env = parent.frame()
)
\end{verbatim}
\end{Usage}
%
\begin{Arguments}
\begin{ldescription}
\item[\code{formula}] A potential formula object to validate (can be any object).

\item[\code{data}] An optional data frame containing the variables.

\item[\code{func\_name}] Character string. Name of the calling function (for error messages).

\item[\code{calling\_env}] The environment in which to look for variables if data is not provided.
Defaults to parent.frame().
\end{ldescription}
\end{Arguments}
%
\begin{Value}
Returns NULL invisibly. Stops with an error if validation fails.
\end{Value}
\HeaderA{validate\_lm2}{Validate Inputs for lm2() Function}{validate.Rul.lm2}
\keyword{internal}{validate\_lm2}
%
\begin{Description}
Validates se\_type and clusters arguments for lm2().
Also handles creating a data frame from vectors if data is not provided.
\end{Description}
%
\begin{Usage}
\begin{verbatim}
validate_lm2(
  formula,
  data = NULL,
  se_type = "HC3",
  se_type_missing = TRUE,
  dots = list(),
  calling_env = parent.frame()
)
\end{verbatim}
\end{Usage}
%
\begin{Arguments}
\begin{ldescription}
\item[\code{formula}] A formula specifying the model.

\item[\code{data}] An optional data frame containing the variables.

\item[\code{se\_type}] The type of standard error to use.

\item[\code{se\_type\_missing}] Logical. Whether se\_type was not explicitly provided by user.

\item[\code{dots}] Additional arguments passed to lm\_robust (to check for clusters).

\item[\code{calling\_env}] The environment in which to look for variables if data is not provided.
\end{ldescription}
\end{Arguments}
%
\begin{Value}
A list containing:
\begin{itemize}

\item{} \code{data}: The data frame to use (either provided or constructed from vectors)
\item{} \code{se\_type}: The validated/adjusted se\_type
\item{} \code{has\_clusters}: Logical indicating if clusters are being used

\end{itemize}

\end{Value}
\HeaderA{validate\_plot}{Validate Inputs for Plotting Functions}{validate.Rul.plot}
\keyword{internal}{validate\_plot}
%
\begin{Description}
Validates inputs for plotting functions that accept either formula syntax
(y \textasciitilde{} group) or standard syntax (y, group), with optional data frame.
\end{Description}
%
\begin{Usage}
\begin{verbatim}
validate_plot(
  y,
  group = NULL,
  data = NULL,
  func_name = "plot",
  require_group = TRUE,
  data_name = NULL
)
\end{verbatim}
\end{Usage}
%
\begin{Arguments}
\begin{ldescription}
\item[\code{y}] A numeric vector, column name, or formula of the form y \textasciitilde{} group.

\item[\code{group}] A vector used to group the data, or a column name if data is provided.
Can be NULL for functions where group is optional.

\item[\code{data}] An optional data frame containing the variables.

\item[\code{func\_name}] Character string. Name of the calling function (for error messages).

\item[\code{require\_group}] Logical. If TRUE, group is required. If FALSE, group can be NULL.

\item[\code{data\_name}] Character string. Name of the data argument (for error messages). 
If NULL, will attempt to infer from the call.
\end{ldescription}
\end{Arguments}
%
\begin{Value}
A list containing:
\begin{itemize}

\item{} \code{y}: Validated y variable (numeric vector)
\item{} \code{group}: Validated group variable (if provided, otherwise NULL)
\item{} \code{y\_name}: Clean variable name for y (for labels)
\item{} \code{group\_name}: Clean variable name for group (for labels)
\item{} \code{y\_name\_raw}: Raw variable name for y (for error messages)
\item{} \code{group\_name\_raw}: Raw variable name for group (for error messages)
\item{} \code{data\_name}: Name of data argument (for error messages)

\end{itemize}

\end{Value}
\HeaderA{validate\_t.test2}{Validate Grouping Variable for t.test2()}{validate.Rul.t.test2}
\keyword{internal}{validate\_t.test2}
%
\begin{Description}
Validates that a grouping variable exists and has exactly 2 levels for t-test.
\end{Description}
%
\begin{Usage}
\begin{verbatim}
validate_t.test2(
  group_var_name,
  data = NULL,
  calling_env = parent.frame(),
  data_name = NULL
)
\end{verbatim}
\end{Usage}
%
\begin{Arguments}
\begin{ldescription}
\item[\code{group\_var\_name}] Character string. Name of the grouping variable (for error messages).

\item[\code{data}] An optional data frame containing the variable.

\item[\code{calling\_env}] The environment in which to look for the variable if data is not provided.

\item[\code{data\_name}] Character string. Name of the data argument (for error messages).
If NULL, will attempt to infer from the call.
\end{ldescription}
\end{Arguments}
%
\begin{Value}
The validated grouping variable (numeric or factor vector). Stops execution with an error message if validation fails.
\end{Value}
\HeaderA{validate\_table2}{Validate Inputs for table2() Function}{validate.Rul.table2}
\keyword{internal}{validate\_table2}
%
\begin{Description}
Validates inputs for table2() function that accepts multiple variables via ...
with optional data frame.
\end{Description}
%
\begin{Usage}
\begin{verbatim}
validate_table2(..., data = NULL, func_name = "table2", data_name = NULL)
\end{verbatim}
\end{Usage}
%
\begin{Arguments}
\begin{ldescription}
\item[\code{...}] One or more variables to be tabulated.

\item[\code{data}] An optional data frame containing the variables.

\item[\code{func\_name}] Character string. Name of the calling function (for error messages).
Default is "table2".

\item[\code{data\_name}] Character string. Name of the data argument (for error messages).
If NULL, will attempt to infer from the call.
\end{ldescription}
\end{Arguments}
%
\begin{Value}
A list containing:
\begin{itemize}

\item{} \code{dots}: List of evaluated variables (ready for base::table)
\item{} \code{dot\_expressions}: List of expressions (for variable name extraction)
\item{} \code{data\_name}: Name of data argument (for error messages)

\end{itemize}

\end{Value}
\HeaderA{var\_labels}{Get or set variable labels}{var.Rul.labels}
\aliasA{var\_labels<\Rdash{}}{var\_labels}{var.Rul.labels<.Rdash.}
%
\begin{Description}
Works like \code{names()}: calling \code{var\_labels(x)} returns labels, and
assignment \code{var\_labels(x) <- value} sets them.
\end{Description}
%
\begin{Usage}
\begin{verbatim}
var_labels(x)

var_labels(x) <- value
\end{verbatim}
\end{Usage}
%
\begin{Arguments}
\begin{ldescription}
\item[\code{x}] A vector or a data frame.
\end{ldescription}
\end{Arguments}
%
\begin{Details}
For a vector, this reads/writes a base-R \code{"label"} attribute.
For a data frame, this reads/writes the \code{"label"} attribute of each
column.
\end{Details}
%
\begin{Value}
- If \code{x} is a data frame: a named character vector with one entry per
column (missing labels are returned as \code{NA\_character\_}).
- Otherwise: a single character string (or \code{NA\_character\_}).
\end{Value}
%
\begin{Examples}
\begin{ExampleCode}
df <- data.frame(x = 1:3, y = 4:6)

# Set labels for all columns
var_labels(df) <- c("this is x", "this is y")
var_labels(df)

# Set a label for a single column
var_labels(df$x) <- "this is x"
var_labels(df$x)

\end{ExampleCode}
\end{Examples}
\printindex{}
\end{document}
